\section{Conclusions}
\label{sec:conclusions}

We built a simulation of the classical and quantum kicked rotor, validated it against exact results, and
used it to follow how a quantum system suppresses chaotic diffusion. The main findings are:
\begin{itemize}
  \item For $K$ above $K_c \approx 0.97$ the classical momentum diffuses, with $D/(K^2/2)$ oscillating around 1
  with period about $2\pi$ in $K$. The Bessel-function expansion reproduces this to an rms miss of 0.04, and
  where it misses, the cause is visible in the measured correlations of the kick sequence. Islands, in
  particular the accelerator modes, make the diffusion anomalous or ensemble dependent.
  \item The quantum energy follows the classical growth at first: exactly after one kick, to order $\hbar^2$
  after two, in detail up to the Ehrenfest time, and statistically up to a break time $t^* \sim D/\hbar^2$.
  Then it saturates at about $0.5$ to $0.8$ times $D^2/\hbar^2$ and only fluctuates.
  \item The time-averaged momentum distribution is exponentially localized, with the shape predicted by
  Izrailev. Its localization length is $\xi \approx 1.2\,D/\hbar^2$ with a scatter of about 20\% over
  $0.2 \le \hbar \le 1$, consistent with the $(1.04 \pm 0.20)\,D/\hbar^2$ from earlier work, and it follows the
  oscillations of the classical $D(K)$.
  \item Floquet theory explains why the growth stops (a discrete quasienergy spectrum makes the energy
  quasi-periodic), and the mapping onto the Anderson model explains why the Floquet states are localized.
  At rational $\hbar/4\pi$ the localization fails, and we reproduced the exact resonance and antiresonance.
  \item Quantum evolution does not amplify errors, while classical chaos amplifies them exponentially, as the
  time-reversal experiment and the comparison of two computers both show.
\end{itemize}

Some observations remain unexplained: the broad bands in $\xi\hbar^2/D$ as $\hbar$ varies, the narrow central
peak of the time-averaged distribution, and the long-range correlations behind the misses of the diffusion
formula at a few values of $K$. Each would make a good follow-up. The code is modular, so further numerical
studies, for example of the statistics of the quasienergies or of how noise destroys localization, can build
on it directly.
